\section{A surface-correspondence criterion for a Mumford square}
\label{sec:mfg-surface-gate}

This section turns the K3 approach to Mumford fourfolds into a precise sufficient condition on one surface correspondence.

Let $(A,\theta)$ be a Hodge-generic Mumford abelian fourfold, let $F$ be
its totally real cubic field, and put $V=H^1(A,\mathbf Q)$.  We use the
split-corestriction case of Mumford's construction, so that
\[
 G^{\mathrm{sc}}_{\mathbf C}\simeq\mathrm{SL}_2^3,
 \qquad V_{\mathbf C}=U_1\otimes U_2\otimes U_3,
 \qquad \dim U_i=2.
\]
Assume $A=\operatorname{Jac}(C)$ for a smooth curve $C$ of genus four.
This hypothesis applies to the two explicit families of
Bouchet--Hanselman--Pieper--Schiavone
\cite[Theorem 1.2]{BHPS}.  The conditional real-multiplication and
Kuga--Satake approach already appears in their Section 5.2; the purpose
here is to state precisely a sufficient surface-correspondence input
and to give the algebraic extraction argument.

\subsection{The two cubic-field representations}

The tensor decompositions are
\begin{align*}
 \operatorname{Sym}^2V_{\mathbf C}
  &=(3,3,3)\oplus(3,1,1)\oplus(1,3,1)\oplus(1,1,3),\\
 \Lambda^2V_{\mathbf C}
  &=(1,1,1)\oplus(1,3,3)\oplus(3,1,3)\oplus(3,3,1).
\end{align*}
Here $(a,b,c)$ records the dimensions of the three irreducible factors.
The adjoint sum in the first line descends to a rational rank-nine
Hodge structure $T$.  The nontrivial sum in the second is
$P=H^2(A,\mathbf Q)_{\mathrm{prim}}$, of rank $27$.
The three factors are permuted transitively by the Galois action defining
$F$, and occur with multiplicity one.  Schur's lemma and descent give
\[
 \operatorname{End}_{\mathrm{HS}}(T)=F,
 \qquad \operatorname{End}_{\mathrm{HS}}(P)=F.
\]
The first equality concerns $T$, and does not assert extra endomorphisms
of $A$: generically $\operatorname{End}^0(A)=\mathbf Q$.
The location of $T$ in $V\otimes V$, rather than in $H^2(A)$, is also
explained in \cite[Section 2.3]{Zhu}.

Identify $V$ with $H^1(C,\mathbf Q)$.  Then
\[
 T\subset \operatorname{Sym}^2V\subset H^2(C\times C,\mathbf Q)
\]
is the unique copy of the rational adjoint representation in $V\otimes V$.
Let $q_0$ be the restriction of the intersection pairing of $C\times C$
to $T$.  This is a K3-type polarization of signature $(2,7)$: the split
real factor contributes Hodge numbers $(1,1,1)$, and the two compact
factors contribute six further $(1,1)$ directions.  There is an
$F$-bilinear ternary form $q_F$ such that
\[
 q_0=\operatorname{Tr}_{F/\mathbf Q}q_F.
\]
The $F$-action is self-adjoint for these forms.

For $P$ use
\[
 Q_2(x,y)=\int_Axy\theta^2.
\]
It identifies $P\otimes P$ with $\operatorname{End}(P)$.  Inside its
Hodge classes, the scalar line is generated by the primitive identity
tensor, which is a divisor-product class.  The complementary plane
\[
 E\simeq F^0:=\ker\operatorname{Tr}_{F/\mathbf Q}
\]
is the usual exceptional plane in $H^4(A^2,\mathbf Q)$; see
\cite[Propositions 4.1--4.2]{Galluzzi}.

\subsection{A nonsquare polarization twist}

Choose a totally positive element
\[
 \lambda\in F^\times_{>0}\setminus\mathbf Q^\times F^{\times2},
 \qquad q_\lambda(x,y)=q_0(\lambda x,y).
\]
Such elements exist.  For example, at a completely split rational prime,
weak approximation gives a totally positive element whose three
valuation parities are $(1,0,0)$, whereas a rational number times a
square has equal parities at those three primes.

The polarized rational Hodge structure $(T,q_\lambda)$ is the
transcendental Hodge structure of a projective K3 surface $S_\lambda$,
up to a positive rational rescaling of the form.  To see this, choose
an even integral lattice of rank nine in the rational quadratic space.
Its discriminant group has length at most nine.  The complementary rank
in the K3 lattice is thirteen, and the complementary signature is
$(1,12)$.  Nikulin's primitive-embedding theorem applies
\cite[Theorem~1.12.4, pp.~122--123]{Nikulin}. Extend the Hodge decomposition by declaring its orthogonal complement to have type $(1,1)$. The K3 period theorem \cite[Theorem~1.7, p.~109]{MorrisonK3} realizes this signed Hodge structure on the K3 lattice.  The complementary lattice
contains a positive class, so the K3 surface is projective.  Since $T$
has no rational $(1,1)$ classes at a Hodge-generic Mumford point, its
transcendental rank is nine and $\rho(S_\lambda)=13$.
A K3 realization of this rank-nine structure with Picard number
thirteen is also given in
\cite[Appendix~A, Theorem~A23]{BBK3Sixfolds}; the argument here
keeps the nonsquare polarization parameter explicit.
A rational rescaling can be absorbed into $\lambda$ without changing
its class modulo $\mathbf Q^\times F^{\times2}$.

\begin{theorem}[One surface correspondence suffices]\label{thm:mumford-surface}
Suppose there is an algebraic cycle
\[
 Z\in\mathrm{CH}^2(S_\lambda\times C\times C)_{\mathbf Q}
\]
whose correspondence action has a nonzero component
\begin{equation}
 \gamma:T(S_\lambda)\longrightarrow
              H^1(C,\mathbf Q)\otimes H^1(C,\mathbf Q).
 \label{eq:mumford-surface-map}
\end{equation}
Then the entire exceptional plane $E\subset H^4(A^2,\mathbf Q)$ is
algebraic.  It is enough to find one effective irreducible surface $Z$
with this property.
\end{theorem}

\begin{proof}
The degree-one K\"unneth projectors of the two curves are algebraic.
The Poincar\'e line bundle and the Abel--Jacobi map give algebraic
inverse identifications of $H^1(C)$ and $H^1(A)$.  Thus the indicated
component is itself an algebraic correspondence into $V\otimes V$.
Since $T$ occurs there with multiplicity one and
$\operatorname{End}_{\mathrm{HS}}(T)=F$, after the fixed Hodge
identification it is multiplication by some $c\in F^\times$ followed
by the inclusion of $T$.

The transcendental projector of $S_\lambda$ is algebraic.  Explicitly,
subtract the degree-zero and degree-four projectors from its diagonal,
and subtract the projector onto its N\'eron--Severi group using the
inverse intersection matrix of a divisor basis.  Its degree-two tensor
is the inverse of $q_\lambda$.  Applying $\gamma$ in both factors gives
an algebraic tensor in
\[
 \operatorname{Sym}^2V\otimes\operatorname{Sym}^2V
        \subset H^4(A^4,\mathbf Q).
\]
Over the three real embeddings of $F$, this is the sum of the three
adjoint Casimir tensors with coefficients
\begin{equation}
 \alpha_i=\sigma_i(c^2/\lambda),\qquad i=1,2,3,
 \label{eq:mumford-casimir-coeff}
\end{equation}
up to a common nonzero rational normalization. Indeed the identification of the adjoint representation with symmetric tensors by the alternating form on $V$ turns $q_0$ into a fixed rational multiple of the trace pairing on $\operatorname{End}(V)$. Each $\mathfrak{sl}_2$ factor acts on $V$ with multiplicity four, so this normalization is the same in all three factors. The inverse polarization
contributes $\lambda^{-1}$ and the two copies of $\gamma$ contribute
$c^2$.  By the choice of $\lambda$, the element $c^2/\lambda$ is not
rational.

Regroup the four $V$ factors and antisymmetrize in the two new pairs.
This is an algebraic operation: permute abelian factors and pull back
along two diagonal maps.  We compute its action on $P$.  Normalize the
$\mathfrak{sl}_2$ Casimir as
\[
 C=H^2/4+(EF+FE)/2,
\]
with eigenvalues $3/4$ on the standard representation and $2$ on its
symmetric square.  Its cross term on $\Lambda^2V$ is
\[
 \Omega_i=C_i/2-\tfrac34\operatorname{id}.
\]
Let $P_j$ be the nine-dimensional summand on which the $j$th factor is
trivial.  Then $\Omega_i$ has eigenvalue $-3/4$ on $P_i$ and $1/4$ on
the other $P_j$.  The recoupled tensor therefore acts on $P_j$ by
\begin{equation}
 \frac{\alpha_1+\alpha_2+\alpha_3}{4}-\alpha_j,
 \label{eq:mumford-casimir-eigen}
\end{equation}
up to one common nonzero factor.  To check that antisymmetrization is
indeed this cross-Casimir map, identify an adjoint element $X$ with the
symmetric form $b_X(u,v)=\varepsilon(Xu,v)$ using the alternating
polarization $\varepsilon$ of $V$.  Antisymmetrizing its Casimir tensor
gives twice
\[
 \sum_X\bigl(b_X(u,w)b_{X^\vee}(v,z)
             -b_X(u,z)b_{X^\vee}(v,w)\bigr),
\]
which is the cross-Casimir pairing on $u\wedge v,w\wedge z$.

Project to $P\otimes P$ and subtract the rational trace.  Both operations
are algebraic, as explained below.  Formula
\eqref{eq:mumford-casimir-eigen} becomes, up to the same nonzero factor,
\[
 -\left(\alpha_j-
      \frac{\alpha_1+\alpha_2+\alpha_3}{3}\right).
\]
This is a nonzero element $a\in F^0$, since $c^2/\lambda\notin\mathbf Q$.
We have therefore constructed one nonzero algebraic exceptional class.

For completeness, no conjectural algebraic projector onto $F$ is used.
Multiplication maps on $A$ give algebraic cohomological K\"unneth
projectors.  Let $L$ be cup product with $\theta$, let
$\mu=\int_A\theta^4/4!$, and use the algebraic lowering operator
\[
 \Lambda(x)=\frac{\theta^3}{6\mu}*x,
\]
where $*$ is Pontryagin product.  On $H^2(A)$ set
\[
 \pi_\theta=\tfrac14L\Lambda,
 \qquad \pi_P=1-\pi_\theta,
 \qquad
 R=\left(\tfrac14\pi_P+\tfrac1{36}\pi_\theta\right)\Lambda^2.
\]
The standard Lefschetz relations give $R=(L^2)^{-1}$ on $H^6(A)$:
$\Lambda^2L^2$ is $4$ on $P$ and $36$ on the polarization line.
Their algebraicity and these constants were proved in Section~\ref{sec:mumford}. Projecting the diagonal and applying $R$
produces an algebraic primitive identity tensor $I_P\in P\otimes P$.
Under $P\otimes P\simeq\operatorname{End}(P)$, subtracting
$\operatorname{Tr}_{\mathbf Q}(a)I_P/27$ is algebraic; the trace is
an intersection pairing with $I_P\theta_1^2\theta_2^2$.

Finally, multiply the tensor representing $a$ by $\theta_1^2$.
It becomes an algebraic degree-zero correspondence on $A$ acting as $a$
on $P$.  Compose it with itself and apply $R$ on the first tensor factor
to obtain an algebraic tensor for $a^2$.  Subtracting the scalar trace
gives
\[
 a^2-\tfrac13\operatorname{Tr}_{F/\mathbf Q}(a^2)\operatorname{id}.
\]
This and $a$ are linearly independent over $\mathbf Q$: otherwise the
nonrational element $a$ of the cubic field $F$ would satisfy a quadratic
polynomial.  They span $F^0$, proving the assertion.  An effective
reducible cycle with a nonzero component has an irreducible component
with that property, which explains the final formulation.
\end{proof}

\subsection{Detection, graphs, and family uniformity}

For an irreducible surface $Z$, the condition in
\eqref{eq:mumford-surface-map} can be tested on a resolution
$\widetilde Z$.  It suffices that, for a nonzero holomorphic two-form
$\omega$ on $S_\lambda$ and some holomorphic one-forms $\eta_1,\eta_2$
on $C$,
\[
 \int_{\widetilde Z}p_S^*\omega\wedge
       \overline{p_1^*\eta_1\wedge p_2^*\eta_2}\ne0.
\]
The form $H^{2,0}(S_\lambda)$ generates its simple rational
transcendental Hodge structure, so this is exactly a nonvanishing test
for the required Hodge component.

\begin{proposition}[The square class of a dominant graph]\label{prop:mumford-graph-twist}
If a dominant rational map $f:C\times C\dashrightarrow S_\lambda$
exists, then $\lambda\in\mathbf Q^\times F^{\times2}$.
\end{proposition}

\begin{proof}
Resolve $f$.  Blowups add only divisor classes.  Pullback on the
transcendental part is therefore a nonzero map into the unique copy of
$T$ in $H^1(C)\otimes H^1(C)$, hence is multiplication by some
$c\in F^\times$.  The projection formula gives
\[
 q_0(cx,cy)=\deg(f)\,q_\lambda(x,y).
\]
Self-adjointness of $F$ and nondegeneracy yield
$c^2=\deg(f)\lambda$, which proves the square-class assertion.
\end{proof}

This proposition does not exclude a general surface correspondence.
Such a correspondence has two projections and need not induce a
similarity of intersection forms.  Nor does it exclude an untwisted
bridge followed by an independently algebraic nonrational RM operator.

If a detecting correspondence as above is constructed over a nonempty
open subset of a fixed algebraic family, all operations in the theorem
are relative algebraic operations after a finite base change.  Their
cohomological coefficients are locally constant, and their cycle
degrees are bounded after spreading out over that open subset.
The resulting two exceptional classes therefore give a familywise
bound for each fixed global class in the exceptional plane; proper
relative Chow parameter spaces extend such a bound by specialization.
Alternatively, a uniform bound on the degrees of exceptional detectors on $A^2$ at an infinite set of points of the Mumford base curve suffices by Corollary~\ref{cor:mumford-family}. A nonzero bridge at a special fibre is not automatically such a detector: its larger Hodge endomorphism algebra need not identify its action with an element of $F$.
An isolated Hodge-generic point with an algebraic correspondence is
only a pointwise input and is not, by itself, a family-uniformity claim.

The theorem leaves precisely the existence of the first detecting
surface unresolved.  The conditional Kuga--Satake route discussed in \cite[Section~5.2, pp.~24--25]{BHPS} likewise requires an algebraic correspondence; an isogeny of the associated abelian varieties does not establish algebraicity of the Kuga--Satake embedding.  The explicit curves
of \cite{BHPS} make the correspondence problem concrete; neither
the preceding argument nor the cited conditional route constructs the
required surface.
